% titel: Blatt 5, A4: Periodizität im PSE (Trends)

\subsection*{Frage}
Wie und warum ändern sich Atomradius, Ionisierungsenergie, Elektronenaffinität
und Metallcharakter innerhalb von Gruppe und Periode? Warum sind die Änderungen
in den Übergangsreihen schwach? Warum zeigen Be und N Maxima der 1.\
Ionisierungsenergie?

\subsection*{Lösung}
\begin{itemize}
  \item \textbf{In der Periode} (z.\ B. Li bis Ne): Kernladung $Z$ wächst,
    Schalenzahl $n$ ist konstant. Durch die höhere Kernladung werden die
    Valenzelektronen stärker angezogen:
    Atomradius $\searrow$, Ionisierungsenergie $\nearrow$, Elektronenaffinität
    $\nearrow$, Metallcharakter $\searrow$.
  \item \textbf{In der Gruppe} (z.\ B. F bis I): Je tiefer, desto mehr Schalen;
    die Abschirmung durch innere Schalen wächst, $Z_{\mathrm{eff}}$ auf das
    Valenzelektron bleibt annähernd gleich. Deshalb: Atomradius $\nearrow$,
    Ionisierungsenergie $\searrow$, Elektronenaffinität $\searrow$,
    Metallcharakter $\nearrow$.
  \item \textbf{Übergangsreihen:} Auffüllen der $(n-1)d$-Orbitale. Die
    hinzukommenden Elektronen liegen im Inneren des Atoms und schirmen die
    zusätzlichen Kernladungen teilweise ab. $Z_{\mathrm{eff}}$ ändert sich daher
    nur langsam $\to$ Radius und Ionisierungsenergie variieren über die Reihe
    nur wenig.
  \item \textbf{Be und N als Maxima:} Be hat die stabile, voll besetzte
    Konfiguration $[\mathrm{He}]\,2s^2$, N die halb besetzte $2p^3$ (Hund'sche
    Stabilität). Zwischen B und C sowie zwischen O und F sinkt die 1.\
    Ionisierungsenergie trotz steigender Kernladung wieder leicht --- genau dort
    gibt je ein Element mit halb- bzw.\ vollbesetzter Unterschale ein Maximum aus.
    Genauer: Be ($2s^2$, voll) steht höher als B, N ($2p^3$, halb) höher als
    O; so entstehen lokale Maxima im sonst aufsteigenden Verlauf der Periode.
\end{itemize}

\begin{ergebnis}
Radius/IE/EA/Metallcharakter folgen $Z_{\mathrm{eff}}$ und $n$: Periode $\to$
Radius kleiner, IE/EA grösser, Metallcharakter kleiner; Gruppe $\to$ umgekehrt.
Übergangsreihen ändern sich schwach wegen d-Abschirmung; Be (volle $2s^2$) und
N (halbe $2p^3$) zeigen IE-Maxima.
\end{ergebnis}

\begin{kontrolle}
Vergleich Pt/K: K bildet Ionen von nur $\mathrm{K^+}$, Pt viel höhere
Ionisierungsenergien --- mit den Regeln erklärbar (Periode 4 vs.\ Übergangsreihe).
Das Modell reproduziert die bekannten Edelgas-Maxima der Ionisierungsenergie.
\end{kontrolle}

\begin{merksatz}
Für Trends immer auf effektive Kernladung und Schalenzahl schauen; halb- und
vollbesetzte Unterschalen => Stabilität => IE-Maximum. In einer Gruppe zählt
die Schalenzahl, in einer Periode die Kernladung.
\end{merksatz}
